Este trabajo presentó un estudio comparativo de representaciones léxicas y basadas en LLM para predecir resultados de contrataciones públicas a partir de documentos de licitación. Los resultados muestran que estos documentos contienen una señal predictiva considerable, pero también que el modelo con mejor ordenamiento en la configuración actual es una línea base léxica clásica: TF--IDF más regresión logística. Al mismo tiempo, el transformer entre fragmentos sobre embeddings de LLM alcanza la mejor calibración y la mayor estabilidad ante cambios del umbral.

La comparación también indica que gran parte de la señal densa útil ya está presente en los embeddings congelados, dado que las líneas base con promedio se mantienen próximas al transformer. Esto sugiere que el modelado explícito de interacciones entre fragmentos no ofrece actualmente una ventaja clara de ordenamiento, aunque todavía puede aportar valor práctico mediante probabilidades más confiables y un comportamiento de decisión más estable. Los trabajos futuros deberían evaluar otros dominios de contratación y otros codificadores congelados.

La principal contribución del estudio es esta comparación empírica e interpretativa. Demuestra que las características léxicas capturan gran parte de la señal predictiva del dominio, que las representaciones densas preentrenadas continúan siendo útiles y que la calidad del ordenamiento, la calibración y la sensibilidad al umbral deben analizarse por separado. Estos hallazgos conectan la evaluación de aprendizaje automático con las necesidades de la gobernanza pública digital: antes de que las herramientas predictivas puedan respaldar la transparencia, la rendición de cuentas o la supervisión ciudadana de las contrataciones, sus señales, limitaciones y comportamiento operativo deben hacerse explícitos.
